\documentclass[11pt]{article}
\usepackage{amsmath,amssymb,amsthm}
\usepackage{algorithm}
\usepackage[noend]{algpseudocode}
\usepackage{fullpage}
\usepackage{hyperref}
\usepackage{enumitem}
\newtheorem{theorem}{Theorem}
\newtheorem{lemma}{Lemma}
\newtheorem{assumption}{Assumption}
\newtheorem{remark}{Remark}
\newcommand{\E}{\mathbb{E}}
\newcommand{\R}{\mathbb{R}}

\begin{document}

\section*{Noisy Gradient Descent (DP‑SGD with Gaussian Noise)}

\section*{Mathematical Formulation}
Let $f:\R^{d}\rightarrow\R$ be the (possibly non‑convex) loss function.  
At iteration $t=0,1,\dots,T-1$ the algorithm maintains a parameter vector $x_t\in\R^{d}$.
\begin{align}
    \text{(Gradient with noise)}\qquad 
    g_t &\;:=\; \nabla f(x_t) \;+\; \xi_t ,                                   \label{eq:gt}\\[4pt]
    \text{(Gaussian noise)}\qquad 
    \xi_t &\sim \mathcal{N}\!\bigl(0,\sigma^{2} I_d\bigr) \;\;\; \text{independently for each $t$,} \label{eq:noise}\\[4pt]
    \text{(Parameter update)}\qquad 
    x_{t+1} &\;=\; x_t \;-\; \eta\, g_t ,                                   \label{eq:update}
\end{align}
where 
\begin{itemize}
    \item $\eta>0$ is the (fixed) learning rate,
    \item $\sigma^{2}>0$ controls the magnitude of the added Gaussian noise (privacy budget),
    \item $I_d$ is the $d\times d$ identity matrix.
\end{itemize}
The algorithm stops after $T$ iterations and returns $x_T$ (or the average $\bar x_T:=\frac1T\sum_{t=0}^{T-1}x_t$).

\section*{Pseudocode}
\begin{algorithm}[H]
\caption{Noisy Gradient Descent (DP‑SGD)}\label{alg:dp-sgd}
\begin{algorithmic}[1]
\Require Initial point $x_0\in\R^d$, learning rate $\eta>0$, noise std.\ $\sigma>0$, number of steps $T$.
\For{$t=0$ \textbf{to} $T-1$}
    \State Compute true gradient $ \nabla f(x_t)$.
    \State Sample $\xi_t \sim \mathcal{N}(0,\sigma^{2}I_d)$.
    \State Form noisy gradient $g_t \gets \nabla f(x_t)+\xi_t$.
    \State Update $x_{t+1} \gets x_t - \eta\, g_t$.
\EndFor
\State \Return $x_T$ (or $\bar x_T=\frac1T\sum_{t=0}^{T-1}x_t$).
\end{algorithmic}
\end{algorithm}

\section*{Dry Run Example}
Consider the one‑dimensional quadratic $f(w)=(w-3)^2$.
\begin{itemize}
    \item True gradient: $\nabla f(w)=2(w-3)$.
    \item Choose $w_0=0$, learning rate $\eta=0.1$, noise variance $\sigma^2=0.01$ (i.e. $\sigma=0.1$).  
    \item For reproducibility we fix a seed and obtain the following noise draws:
          $\xi_0=0.05$, $\xi_1=-0.08$, $\xi_2=0.02$.
\end{itemize}
\[
\begin{array}{c|c|c|c|c}
t & w_t & \nabla f(w_t) & g_t=\nabla f(w_t)+\xi_t & f(w_{t+1})\\ \hline
0 & 0 & -6 & -6+0.05=-5.95 & f\bigl(0-0.1\cdot(-5.95)\bigr)=f(0.595)= (0.595-3)^2=5.78\\
1 & 0.595 & 2(0.595-3)=-4.81 & -4.81-0.08=-4.89 & f\bigl(0.595-0.1\cdot(-4.89)\bigr)=f(1.084)= (1.084-3)^2=3.68\\
2 & 1.084 & 2(1.084-3)=-3.832 & -3.832+0.02=-3.812 & f\bigl(1.084-0.1\cdot(-3.812)\bigr)=f(1.465)= (1.465-3)^2=2.36
\end{array}
\]
After three iterations the objective has decreased from $f(w_0)=9$ to $f(w_3)\approx2.36$.

\newpage

\section*{Assumptions}
\begin{assumption}[Smoothness]\label{as:smooth}
$f$ is $L$‑smooth, i.e.
\[
\forall x,y\in\R^d,\qquad 
\|\nabla f(x)-\nabla f(y)\|\le L\|x-y\|.
\]
Equivalently,
\[
f(y)\le f(x)+\langle \nabla f(x),y-x\rangle+\frac{L}{2}\|y-x\|^{2}.
\]
\end{assumption}
\begin{assumption}[Bounded Below]\label{as:lower}
There exists $f_{\inf}>-\infty$ such that $f(x)\ge f_{\inf}$ for all $x$.
\end{assumption}
\begin{assumption}[Noise Distribution]\label{as:noise}
The added noise $\{\xi_t\}_{t\ge0}$ satisfies
\[
\xi_t\sim\mathcal N(0,\sigma^{2}I_d),\qquad 
\E[\xi_t]=0,\qquad 
\E\bigl[\|\xi_t\|^{2}\bigr]=\sigma^{2}d.
\]
Moreover the noise is independent of $x_t$ and of all previous randomness.
\end{assumption}
\begin{assumption}[Learning Rate]\label{as:lr}
The stepsize obeys $0<\eta\le \frac{1}{L}$.
\end{assumption}

\section*{Main Convergence Theorem}
\begin{theorem}[Expected Stationarity]\label{thm:main}
Under Assumptions \ref{as:smooth}–\ref{as:lr}, after $T$ iterations of the algorithm,
\[
\frac{1}{T}\sum_{t=0}^{T-1}\E\!\bigl[\|\nabla f(x_t)\|^{2}\bigr]
\;\le\;
\frac{2\bigl(f(x_0)-f_{\inf}\bigr)}{\eta T}
\;+\;
L\eta\,\sigma^{2}d .
\]
Consequently, choosing $\eta =\frac{c}{\sqrt{T}}$ with $c\le \frac{1}{L}$ yields
\[
\frac{1}{T}\sum_{t=0}^{T-1}\E\!\bigl[\|\nabla f(x_t)\|^{2}\bigr]
\;=\; O\!\left(\frac{1}{\sqrt{T}}\right) \;+\; O\!\left(\frac{\sigma^{2}d}{\sqrt{T}}\right).
\]
\end{theorem}

\section*{Theoretical Properties}
\begin{itemize}
    \item \textbf{Stability.} Because the update uses a stepsize $\eta\le1/L$, the iterates remain bounded in expectation; the additive Gaussian term does not destabilise the method.
    \item \textbf{Hyper‑parameter sensitivity.} The bound contains the product $L\eta\sigma^{2}d$; larger noise $\sigma$ or larger stepsize $\eta$ increase the stationary‑point error linearly. The optimal $\eta$ balances the $1/(\eta T)$ term and the $L\eta\sigma^{2}d$ term, giving $\eta\asymp 1/\sqrt{T}$.
    \item \textbf{Comparison to vanilla SGD.} Setting $\sigma=0$ recovers the standard non‑convex SGD bound $\frac{2(f(x_0)-f_{\inf})}{\eta T}$, i.e. $O(1/T)$ for constant $\eta$. The privacy‑induced noise adds an irreducible $O(\sigma^{2}d/\sqrt{T})$ term.
\end{itemize}

\section*{Discussion}
The theorem shows that even with Gaussian noise sufficient for differential privacy, the algorithm converges to an \emph{approximate stationary point}. The rate matches the best known for DP‑SGD under smooth non‑convex objectives. By decreasing $\eta$ as $1/\sqrt{T}$ the privacy‑induced error decays at the same order as the optimisation error, yielding a meaningful trade‑off between privacy (through $\sigma$) and utility.

\newpage

\section*{Preliminaries (Definitions and Lemmas)}
\begin{lemma}[Smoothness inequality]\label{lem:smooth}
If $f$ is $L$‑smooth then for any $x,y\in\R^{d}$,
\[
f(y)\le f(x)+\langle \nabla f(x),y-x\rangle+\frac{L}{2}\|y-x\|^{2}.
\]
\end{lemma}
\begin{proof}
Standard consequence of the $L$‑Lipschitz gradient; see e.g. \cite{Nesterov2004}.
\end{proof}

\begin{lemma}[Unbiasedness of the noisy gradient]\label{lem:unbiased}
For each $t$, $\E[g_t\mid x_t]=\nabla f(x_t)$.
\end{lemma}
\begin{proof}
From definition \eqref{eq:gt} and Assumption~\ref{as:noise},
\[
\E[g_t\mid x_t]=\nabla f(x_t)+\E[\xi_t]=\nabla f(x_t).
\]
\end{proof}

\section*{Main Proof (Step‑by‑Step Derivation)}
We prove Theorem~\ref{thm:main}. Throughout the proof all expectations are taken with respect to the randomness up to iteration $t$ unless otherwise noted.

\noindent\textbf{Step 1 (Smoothness upper bound).}  
Using Lemma~\ref{lem:smooth} with $x=x_t$, $y=x_{t+1}=x_t-\eta g_t$,
\[
f(x_{t+1})\le f(x_t)-\eta\langle \nabla f(x_t),g_t\rangle
+\frac{L\eta^{2}}{2}\|g_t\|^{2}.                                 \tag{1}
\]

\noindent\textbf{Step 2 (Expand the inner product).}  
Insert $g_t=\nabla f(x_t)+\xi_t$,
\[
\langle \nabla f(x_t),g_t\rangle
= \|\nabla f(x_t)\|^{2}+\langle \nabla f(x_t),\xi_t\rangle .      \tag{2}
\]

\noindent\textbf{Step 3 (Expand the squared norm).}  
\[
\|g_t\|^{2}
= \|\nabla f(x_t)\|^{2}+2\langle \nabla f(x_t),\xi_t\rangle+\|\xi_t\|^{2}. \tag{3}
\]

\noindent\textbf{Step 4 (Plug (2) and (3) into (1)).}  
\begin{align}
f(x_{t+1})
&\le f(x_t)-\eta\Bigl(\|\nabla f(x_t)\|^{2}+\langle \nabla f(x_t),\xi_t\rangle\Bigr)  \nonumber\\
&\quad +\frac{L\eta^{2}}{2}\Bigl(\|\nabla f(x_t)\|^{2}+2\langle \nabla f(x_t),\xi_t\rangle+\|\xi_t\|^{2}\Bigr). \tag{4}
\end{align}

\noindent\textbf{Step 5 (Take conditional expectation).}  
Condition on $x_t$ and use $\E[\xi_t]=0$, $\E[\langle \nabla f(x_t),\xi_t\rangle]=0$,
\[
\E\!\bigl[f(x_{t+1})\mid x_t\bigr]
\le f(x_t)-\eta\|\nabla f(x_t)\|^{2}
+\frac{L\eta^{2}}{2}\Bigl(\|\nabla f(x_t)\|^{2}+\E[\|\xi_t\|^{2}]\Bigr). \tag{5}
\]

\noindent\textbf{Step 6 (Insert noise variance).}  
By Assumption~\ref{as:noise}, $\E[\|\xi_t\|^{2}]=\sigma^{2}d$. Hence
\[
\E\!\bigl[f(x_{t+1})\mid x_t\bigr]
\le f(x_t)-\Bigl(\eta-\frac{L\eta^{2}}{2}\Bigr)\|\nabla f(x_t)\|^{2}
+\frac{L\eta^{2}}{2}\sigma^{2}d .                                   \tag{6}
\]

\noindent\textbf{Step 7 (Use the stepsize bound).}  
Assumption~\ref{as:lr} gives $\eta\le 1/L$, which implies
\[
\eta-\frac{L\eta^{2}}{2}\;\ge\;\frac{\eta}{2}.                     \tag{7}
\]

\noindent\textbf{Step 8 (Simplify (6) with (7)).}  
\[
\E\!\bigl[f(x_{t+1})\mid x_t\bigr]
\le f(x_t)-\frac{\eta}{2}\|\nabla f(x_t)\|^{2}
+\frac{L\eta^{2}}{2}\sigma^{2}d .                                   \tag{8}
\]

\noindent\textbf{Step 9 (Take full expectation).}  
Denote $\mathcal{F}_t$ the sigma‑algebra generated up to iteration $t$.
Taking expectation of (8) yields
\[
\E\!\bigl[f(x_{t+1})\bigr]
\le \E\!\bigl[f(x_t)\bigr]-\frac{\eta}{2}\E\!\bigl[\|\nabla f(x_t)\|^{2}\bigr]
+\frac{L\eta^{2}}{2}\sigma^{2}d .                                   \tag{9}
\]

\noindent\textbf{Step 10 (Sum over $t=0,\dots,T-1$).}  
Telescoping the left‑hand side,
\[
\E\!\bigl[f(x_T)\bigr]-f(x_0)
\le -\frac{\eta}{2}\sum_{t=0}^{T-1}\E\!\bigl[\|\nabla f(x_t)\|^{2}\bigr]
+ \frac{L\eta^{2}}{2}\sigma^{2}d\,T .                               \tag{10}
\]

\noindent\textbf{Step 11 (Use lower bound on $f$).}  
Since $f(x_T)\ge f_{\inf}$,
\[
f_{\inf}-f(x_0)
\le -\frac{\eta}{2}\sum_{t=0}^{T-1}\E\!\bigl[\|\nabla f(x_t)\|^{2}\bigr]
+ \frac{L\eta^{2}}{2}\sigma^{2}d\,T .                               \tag{11}
\]

\noindent\textbf{Step 12 (Rearrange).}  
\[
\sum_{t=0}^{T-1}\E\!\bigl[\|\nabla f(x_t)\|^{2}\bigr]
\le \frac{2\bigl(f(x_0)-f_{\inf}\bigr)}{\eta}
+ L\eta\sigma^{2}d\,T .                                            \tag{12}
\]

\noindent\textbf{Step 13 (Divide by $T$).}  
\[
\frac{1}{T}\sum_{t=0}^{T-1}\E\!\bigl[\|\nabla f(x_t)\|^{2}\bigr]
\le \frac{2\bigl(f(x_0)-f_{\inf}\bigr)}{\eta T}
+ L\eta\sigma^{2}d .                                                \tag{13}
\]

Equation \eqref{13} is exactly the bound claimed in Theorem~\ref{thm:main}. ∎

\section*{Rate Derivation (With Constants)}
Choose $\eta = \frac{c}{\sqrt{T}}$ where $c\le 1/L$ to satisfy Assumption~\ref{as:lr}. Plugging into \eqref{13} gives
\[
\frac{1}{T}\sum_{t=0}^{T-1}\E\!\bigl[\|\nabla f(x_t)\|^{2}\bigr]
\le \frac{2\bigl(f(x_0)-f_{\inf}\bigr)}{c}\,\frac{1}{\sqrt{T}}
+ Lc\sigma^{2}d\,\frac{1}{\sqrt{T}} .
\]
Hence
\[
\boxed{\;
\frac{1}{T}\sum_{t=0}^{T-1}\E\!\bigl[\|\nabla f(x_t)\|^{2}\bigr]
= O\!\Bigl(\frac{1}{\sqrt{T}}\Bigr)
\;+\;
O\!\Bigl(\frac{\sigma^{2}d}{\sqrt{T}}\Bigr)
\;}
\]
which proves the stated $O(1/\sqrt{T})$ convergence rate.

\section*{Special Cases}
\begin{enumerate}[label=\arabic*)]
    \item \textbf{No privacy noise ($\sigma=0$).}  
    Bound \eqref{13} reduces to $\frac{2(f(x_0)-f_{\inf})}{\eta T}$, i.e. the classic non‑convex SGD rate $O(1/T)$ for constant $\eta$.
    \item \textbf{Convex smooth case.}  
    If $f$ is additionally convex, the same analysis yields a bound on suboptimality $f(\bar x_T)-f^\star$ of order $O(1/\sqrt{T})$ plus the privacy term.
    \item \textbf{Large batch / reduced noise.}  
    In practice one may average gradients over a minibatch of size $B$, which divides the noise variance by $B$; the term $L\eta\sigma^{2}d$ is replaced by $L\eta\sigma^{2}d/B$.
\end{enumerate}

\begin{thebibliography}{9}
\bibitem{Nesterov2004}
Y.~Nesterov, \emph{Introductory Lectures on Convex Optimization: A Basic Course}, Kluwer Academic Publishers, 2004.
\end{thebibliography}

\end{document}